\section{Methods}
\label{app:methods}

\paragraph*{Metrics.}
The $2\times2$ blocks $p$ of $P=I_2\otimes p$, acting on $(x_i,v_i)$, are, rounded to four digits (the exact values are in the repository),
\begin{equation*}
\dtwo:\ \begin{pmatrix}1&0\\0&2\end{pmatrix},
\end{equation*}
\begin{equation*}
\optc{1.25}:\ \begin{pmatrix}-0.5385&-1.0770\\-0.8616&0.4308\end{pmatrix},
\end{equation*}
\begin{equation*}
\optc{5}:\ \begin{pmatrix}-2.0100&-4.0200\\-0.8040&0.4020\end{pmatrix},
\end{equation*}
and $p=I$ for $P=I$.
They have condition numbers 2, 1.25 and 5, $\kP=0$, $0$ and $-6$, and unit position variance of $p^{-1}p^{-\top}$.

\paragraph*{Training details.}
Batches of 1024 samples are drawn from a fixed training set of $200\,000$ samples of the mixture; half of each batch has $\log\sigma$ uniform in $[\log\smin,\log\smax]$, half in $[\log\smin,\log 10\smin]$.
The learning rate rises linearly to $2\times10^{-3}$ over the first 1000 steps and then follows a cosine decay to $10^{-5}$ (Adam, $\beta_1=0.9$, $\beta_2=0.999$); the exponential moving average has decay $0.9995$.
The noise level enters through Fourier features of $u=\ln(\sigma/\smin)/\ln(\smax/\smin)$.
No validation set is used; the score error of Sec.~\ref{sec:learning} is measured on fresh samples of $p_t$.

\paragraph*{Grids and coupling.}
The finest grid has $K=4496$ steps ($h\approx0.0020$).
The coarser grids used for the SDE have $K=281$, $562$, $1124$ ($h\approx0.032$, $0.016$, $0.008$).
Their Gaussian increments are normalized sums of 16, 8 and 4 consecutive increments of the finest grid, so all levels share one Brownian path.
The grid with $K=140$ ($h\approx0.064$) uses normalized sums of 32 draws from the same stream; its time points do not line up with the finest grid, so it is not strictly coupled to the other levels.
For sampling seed $s$, the prior draw uses seed $100+s$ and the noise seed $200+s$; every metric and network uses the same seeds.
The prior is sampled exactly from $p_T$.

\paragraph*{Windows and perturbations.}
For the window experiments of Sec.~\ref{sec:mechanism} the reverse time is split into eight bands at $\sigma=40$, $1$, $0.3$, $0.1$, $0.05$, $0.03$, $0.02$, $0.01$, $0.005$.
A step belongs to the band that contains $\sigma(t_k)$ at its start.
The perturbations of Sec.~\ref{sec:h9} are the learned-minus-exact path differences, in $Z$, of the final outliers of the twelve broken runs with $P=I$, taken at the injection level.
They are added to all points of the prior draw (seed 100) whose Mahalanobis distance to the nearest component of $p_T$ is between 3.5 and 5 (291 points), four perturbations per point.
A run whose final state is non-finite counts as $A=\infty$ and enters the percentiles as the largest value; nothing is clipped. In all 16 networks the 90th percentiles of $A$ were finite.

\paragraph*{Statistics.}
Breakage counts are compared with one-sided Fisher exact tests~\cite{fisher1922}, with Holm's correction where two secondary comparisons share a family~\cite{holm1979}.
Two runs of the same network share the network and are not independent, so the network is the unit of the primary tests.
Amplification ratios are compared with a one-sided Mann--Whitney test~\cite{mannwhitney1947}.

\paragraph*{Pre-registration.}
Each stage was pre-registered in a time-stamped file before its runs, with the criteria of every test (Table~\ref{tab:prereg}).
Changes made after registration were recorded with their time and reason.
One change concerned an unregistered exploratory threshold before the runs (stage B9).
In stage B11, a small code check ($N=400$) displayed an amplification factor at a level that was not used for any test; this was recorded.
The prediction of Test~\ref{test:h7} was written and hashed before the new networks were sampled.

\paragraph*{Networks.}
Table~\ref{tab:nets} lists every network with its score error and whether its probability-flow ODE broke.
\begin{table}[h]
\caption{All 28 networks. $e_Z$: relative mean-square score error near $\smin$ on typical points (tail: Mahalanobis distance 3.5--5). Broken: probability-flow ODE, either seed. Max $|L|$ and outliers over the two runs.}
\label{tab:nets}
\footnotesize
\begin{ruledtabular}
\begin{tabular}{lccccc}
network & $e_Z$ & tail & broken & max $|L|$ & outliers \\
\hline
$P=I$ 1 & 0.118 & 0.122 & yes & 90 & 6 \\
2 & 0.117 & 0.100 & yes & $\infty$ & 14 \\
3 & 0.112 & 0.105 & yes & $\infty$ & 11 \\
4 & 0.126 & 0.097 & yes & 44 & 18 \\
5 & 0.132 & 0.094 & no & 0.273 & 4 \\
6 & 0.129 & 0.090 & no & 0.720 & 3 \\
7 & 0.113 & 0.086 & yes & $\infty$ & 8 \\
8 & 0.115 & 0.096 & yes & $\infty$ & 15 \\
9 & 0.112 & 0.087 & yes & 117 & 11 \\
10 & 0.130 & 0.095 & yes & $\infty$ & 16 \\
11 & 0.137 & 0.091 & yes & 43 & 6 \\
12 & 0.116 & 0.088 & yes & $\infty$ & 24 \\
$\dtwo$ 1 & 0.195 & 0.120 & no & 0.642 & 3 \\
2 & 0.196 & 0.119 & no & 0.179 & 7 \\
3 & 0.190 & 0.142 & yes & 3.242 & 6 \\
4 & 0.191 & 0.115 & no & 0.110 & 2 \\
5 & 0.220 & 0.139 & no & 0.093 & 3 \\
6 & 0.184 & 0.113 & no & 0.291 & 9 \\
7 & 0.202 & 0.114 & yes & 4.912 & 8 \\
8 & 0.189 & 0.106 & no & 0.187 & 3 \\
$\optc{1.25}$ 1 & 0.087 & 0.121 & no & 0.071 & 1 \\
2 & 0.092 & 0.141 & no & 0.074 & 0 \\
3 & 0.105 & 0.152 & no & 0.099 & 3 \\
4 & 0.089 & 0.115 & no & 0.068 & 1 \\
$\optc{5}$ 1 & 0.298 & 0.404 & yes & $\infty$ & 5 \\
2 & 0.258 & 0.456 & no & 0.536 & 6 \\
3 & 0.327 & 0.519 & no & 0.408 & 6 \\
4 & 0.325 & 0.303 & no & 0.077 & 3 \\
\end{tabular}
\end{ruledtabular}
\end{table}

